\clearpage
\VerdierPhysicalPageAnchor{0026}
多重加法的関手の複体への自然な拡張は、シフトと必ずしも可換ではないが、
「同型を除いて」可換であり、その同型は実際には符号を除いて定まる。
「シフト関手」の合成についてこれらの同型を整合的に選べず、そのため符号
が現れる。第 1 節ではこの符号を研究する（付録も参照）。第 2 節では多重
複体、単純複体、「付随単純複体」関手、ホモトピーという古典的概念を定義
し、第 3 節では複体の射の錐（「写像錐」または「写像柱」）を研究する。
これは以後の議論に不可欠な構成である。

\VerdierSemanticAnchor{intro-6-plan-derived-category}
アーベル圏 $\mathcal A$ の導来圏 $\mathbf D(\mathcal A)$ は二段階で構成
される。第一段階で、$\mathcal A$ の複体をホモトピーを除いて考えた圏
$\mathbf K(\mathcal A)$ を導入する。次に $\mathbf K(\mathcal A)$ の
擬同型を形式的に可逆にして $\mathbf D(\mathcal A)$ を得る。圏
$\mathbf K(\mathcal A)$ と $\mathbf D(\mathcal A)$ は必ずしもアーベル
圏ではないが、錐の構成から定まる図式の族、すなわち\emph{完全三角}という
付加構造をもつ。これらの圏では、完全三角がアーベル圏の完全列の役割を
果たす。

\VerdierSemanticAnchor{intro-6-plan-chapter-two}
Puppe [2] によって見いだされ研究されたこの付加構造をもつ加法圏を
\emph{三角圏}と呼び、第 II 章でわれわれの観点から研究する。第 1 節は
三角圏の初等的結果に充てる。第 2 節では三角圏の一般的枠組みで、
$\mathbf K(\mathcal A)$ の擬同型を形式的に可逆にする問題が「分数計算」
[16] によって簡潔に解かれることを示す。第 3 節では Freyd [17] の定理を
証明する：三角圏 $\mathcal D$ は普遍的にアーベル圏
$\mathrm A(\mathcal D)$ へ埋め込まれる。この圏は、同時に入射的かつ射影的
な対象の族 $E$ をもち、$\mathrm A(\mathcal D)$ の各対象は $E$ のある
対象の部分対象にも商対象にも同型である。$E$ が定める
$\mathrm A(\mathcal D)$ の充満部分圏は $\mathcal D$ と同型である。最後
に第 4 節では、三角圏またはアーベル圏に値をとる「スペクトル対象」を研究
する。ここでは [1] の定義をわれわれの状況に合わせるだけである。

\VerdierSemanticAnchor{intro-6-plan-chapter-three}
第 III 章は導来圏そのものの研究に充てる。$X,Y$ をアーベル圏
$\mathcal A$ の導来圏 $\mathbf D(\mathcal A)$ の二対象とする。
$\mathcal A$ が十分な入射対象をもち、$Y$ の成分が十分小さい次数で零に
なるならば、$X$ から $Y$ への次数 $n$ の射は
